\documentclass[11pt,a4paper]{article}
% preamble.tex — estilo común de todos los apuntes Froth.
% Cada apunte hace % preamble.tex — estilo común de todos los apuntes Froth.
% Cada apunte hace % preamble.tex — estilo común de todos los apuntes Froth.
% Cada apunte hace \input{preamble} justo después de \documentclass.
\usepackage[margin=2.4cm]{geometry}
\usepackage{xcolor}
\definecolor{litblue}{RGB}{31,79,121}
\definecolor{litgreen}{RGB}{27,94,32}
\definecolor{litorange}{RGB}{230,81,0}
\usepackage{parskip}
\usepackage{titlesec}
\titleformat{\section}{\Large\bfseries\color{litblue}}{\thesection}{0.6em}{}
\titleformat{\subsection}{\large\bfseries\color{litblue}}{\thesubsection}{0.6em}{}
\usepackage{tcolorbox}
\usepackage{fancyhdr}
\pagestyle{fancy}
\fancyhf{}
\fancyhead[L]{\small\color{litblue}Froth — Apuntes de ML}
\fancyhead[R]{\small\thepage}
\renewcommand{\headrulewidth}{0.4pt}
\usepackage[colorlinks=true,linkcolor=litblue,urlcolor=litblue]{hyperref}

% Cabecera de cada apunte: \cabecera{numero}{titulo}
\newcommand{\cabecera}[2]{%
  \begin{tcolorbox}[colback=litblue,coltext=white,colframe=litblue,arc=2mm]
    {\Large\bfseries Apunte #1 — #2}\\[3pt]
    {\small Proyecto Froth · aprender Machine Learning construyendo}
  \end{tcolorbox}\vspace{0.8em}}

% Cajas temáticas
\newtcolorbox{concepto}[1]{colback=blue!4,colframe=litblue,title={Concepto: #1},arc=1.5mm}
\newtcolorbox{practica}{colback=green!5,colframe=litgreen,title={Pruébalo tú (teclado en mano)},arc=1.5mm}
\newtcolorbox{ojo}{colback=orange!8,colframe=litorange,title={Ojo / error típico},arc=1.5mm}
\newtcolorbox{resumen}{colback=gray!8,colframe=gray!60,title={En una frase},arc=1.5mm}

\newcommand{\codigo}[1]{\texttt{#1}}
 justo después de \documentclass.
\usepackage[margin=2.4cm]{geometry}
\usepackage{xcolor}
\definecolor{litblue}{RGB}{31,79,121}
\definecolor{litgreen}{RGB}{27,94,32}
\definecolor{litorange}{RGB}{230,81,0}
\usepackage{parskip}
\usepackage{titlesec}
\titleformat{\section}{\Large\bfseries\color{litblue}}{\thesection}{0.6em}{}
\titleformat{\subsection}{\large\bfseries\color{litblue}}{\thesubsection}{0.6em}{}
\usepackage{tcolorbox}
\usepackage{fancyhdr}
\pagestyle{fancy}
\fancyhf{}
\fancyhead[L]{\small\color{litblue}Froth — Apuntes de ML}
\fancyhead[R]{\small\thepage}
\renewcommand{\headrulewidth}{0.4pt}
\usepackage[colorlinks=true,linkcolor=litblue,urlcolor=litblue]{hyperref}

% Cabecera de cada apunte: \cabecera{numero}{titulo}
\newcommand{\cabecera}[2]{%
  \begin{tcolorbox}[colback=litblue,coltext=white,colframe=litblue,arc=2mm]
    {\Large\bfseries Apunte #1 — #2}\\[3pt]
    {\small Proyecto Froth · aprender Machine Learning construyendo}
  \end{tcolorbox}\vspace{0.8em}}

% Cajas temáticas
\newtcolorbox{concepto}[1]{colback=blue!4,colframe=litblue,title={Concepto: #1},arc=1.5mm}
\newtcolorbox{practica}{colback=green!5,colframe=litgreen,title={Pruébalo tú (teclado en mano)},arc=1.5mm}
\newtcolorbox{ojo}{colback=orange!8,colframe=litorange,title={Ojo / error típico},arc=1.5mm}
\newtcolorbox{resumen}{colback=gray!8,colframe=gray!60,title={En una frase},arc=1.5mm}

\newcommand{\codigo}[1]{\texttt{#1}}
 justo después de \documentclass.
\usepackage[margin=2.4cm]{geometry}
\usepackage{xcolor}
\definecolor{litblue}{RGB}{31,79,121}
\definecolor{litgreen}{RGB}{27,94,32}
\definecolor{litorange}{RGB}{230,81,0}
\usepackage{parskip}
\usepackage{titlesec}
\titleformat{\section}{\Large\bfseries\color{litblue}}{\thesection}{0.6em}{}
\titleformat{\subsection}{\large\bfseries\color{litblue}}{\thesubsection}{0.6em}{}
\usepackage{tcolorbox}
\usepackage{fancyhdr}
\pagestyle{fancy}
\fancyhf{}
\fancyhead[L]{\small\color{litblue}Froth — Apuntes de ML}
\fancyhead[R]{\small\thepage}
\renewcommand{\headrulewidth}{0.4pt}
\usepackage[colorlinks=true,linkcolor=litblue,urlcolor=litblue]{hyperref}

% Cabecera de cada apunte: \cabecera{numero}{titulo}
\newcommand{\cabecera}[2]{%
  \begin{tcolorbox}[colback=litblue,coltext=white,colframe=litblue,arc=2mm]
    {\Large\bfseries Apunte #1 — #2}\\[3pt]
    {\small Proyecto Froth · aprender Machine Learning construyendo}
  \end{tcolorbox}\vspace{0.8em}}

% Cajas temáticas
\newtcolorbox{concepto}[1]{colback=blue!4,colframe=litblue,title={Concepto: #1},arc=1.5mm}
\newtcolorbox{practica}{colback=green!5,colframe=litgreen,title={Pruébalo tú (teclado en mano)},arc=1.5mm}
\newtcolorbox{ojo}{colback=orange!8,colframe=litorange,title={Ojo / error típico},arc=1.5mm}
\newtcolorbox{resumen}{colback=gray!8,colframe=gray!60,title={En una frase},arc=1.5mm}

\newcommand{\codigo}[1]{\texttt{#1}}

\begin{document}
\cabecera{40}{Títulos de cluster con sentido: KeyBERT (no oraciones recortadas)}

\section{El problema (lo viste tú)}

La v1 de los subtítulos tomaba la \emph{oración} más central del cluster y la cortaba a
12 palabras. Resultado: encabezados como ``Considering both\ldots'', ``In the present
st\ldots'' --- fragmentos que insinúan el tema pero se sienten recortados, como una
biblia a medio leer. Y las keywords viejas (c-TF-IDF: ``dpa, dda, lepidolite'') tampoco:
clasifican bien pero \emph{no comunican}.

\section{Cómo lo resuelve el campo: KeyBERT}

El estándar (KeyBERT, y la capa de representación de BERTopic) no recorta oraciones:
extrae \textbf{keyphrases} cortas y las rankea por SIGNIFICADO.

\begin{concepto}{KeyBERTInspired en 3 pasos}
\textbf{1. Candidatos}: n-gramas de 2--3 palabras del texto del cluster (con el mismo
vectorizador endurecido de \texttt{cluster.py}: sin dígitos, elementos químicos
salvados). \\
\textbf{2. Rankear}: embeber cada candidato con SPECTER y medir su coseno al
\emph{centroide} del cluster. \\
\textbf{3. Diversificar}: MMR para que las frases elegidas no sean casi iguales.
\end{concepto}

Lo elegante: son las MISMAS piezas que Froth ya tenía (SPECTER + coseno + centroide +
MMR de la Fase 7). Por eso se hizo a mano, sin dependencia nueva --- el patrón del
proyecto: construir cada pieza para entenderla.

\section{Dos tropiezos que enseñan (tus datos reales)}

El primer intento sacó títulos con ruido de OCR (``limu'', ``nmic'', ``cgfs'') y
genéricos (``ore flotation'' donde debía decir ``lepidolite flotation''). Dos arreglos:

\begin{itemize}
  \item \textbf{\texttt{min\_df} $\geq$ 2}: exigir que la frase aparezca en VARIOS papers
        del cluster. La basura de OCR vive en un solo abstract $\to$ fuera.
  \item \textbf{Score contrastivo} (la idea de c-TF-IDF, pero en el espacio de
        embeddings):
        \[
        \mathrm{score} = \cos(\mathrm{frase}, \mathrm{centroide}_c)
                       - 0.4 \, \cos(\mathrm{frase}, \mathrm{centroide}_{\mathrm{global}})
        \]
        Premia lo que está cerca de ESTE cluster y lejos del corpus entero. Así
        ``lepidolite flotation'' (específico) le gana a ``ore flotation'' (genérico,
        cerca del centro de todo).
\end{itemize}

\begin{ojo}
El keyphrase más cercano al centroide NO es el mejor título: suele ser el más genérico.
La especificidad (lejos del promedio global) es lo que hace que un título \emph{diga}
algo. Misma lección que el h-index y el filtro de relevancia: la señal está en el
contraste, no en el valor absoluto.
\end{ojo}

\begin{practica}
El título corto va al encabezado; la oración centroide completa se guarda en la columna
\texttt{summary} y aparece en el \textbf{hover} --- corto arriba, detalle al pasar el
mouse, justo como lo pediste.
\end{practica}

\begin{resumen}
Títulos de cluster = keyphrases de 2--3 palabras extraídas por KeyBERT (candidatos +
SPECTER + coseno al centroide + MMR), con \texttt{min\_df} contra el ruido y un score
contrastivo (cluster $-$ global) para preferir lo específico. Cortos, con sentido, sin
recortar oraciones; la frase completa vive en el hover.
\end{resumen}

\end{document}
